\section{从数学到更广义的科学发现}

AI for Math 之所以重要，还因为数学是科学发现中最形式化、最可验证的一块。如果 AI 能在数学中形成“提出猜想--构造证明--机器验证--人类评价”的闭环，那么类似方法可能迁移到物理、材料、程序验证和算法发现。

\subsection{为什么数学是试验场}

数学有三个优势。第一，正确性可以被形式化系统强验证。第二，历史数据和定理库具有结构化积累。第三，数学家社区对优雅、解释力和重要性有清晰但难以量化的评价传统。这使它既适合机器训练，也能暴露机器训练的不足。

\begin{knowledgebox}{科学发现的三类反馈}
数学提供形式验证；代码提供测试和运行反馈；实验科学提供观测和实验反馈。三者都是 AI 从语言生成走向真实发现的关键环境。EP139 的 Agent、EP138 的模型实验室、EP137 的 AI for Math，其实都在寻找可反馈环境。
\end{knowledgebox}

\subsection{人机协作的终局}

AI for Math 的理想终局未必是“替代数学家”。更现实也更有价值的形态，是 AI 帮数学家形式化繁琐证明、搜索 lemma、发现反例、整理文献、提出候选猜想；数学家则负责问题品味、方向选择、解释和最终判断。这样，AI 成为研究伙伴，而不是单纯自动答题器。

\begin{importantbox}{从工具到伙伴}
当 AI 只做局部步骤，它是工具；当 AI 能理解上下文、积累经验、提出候选方向并接受验证，它开始像研究伙伴。AI for Math 是这个转变的高难度样板。
\end{importantbox}

\subsection{本章小结}

数学是 AI 进入科学发现的高价值试验场。它的强验证性让训练更清晰，它的审美和直觉又提醒我们：科学发现不能被简化成单一指标。

